\subsection{다운스트림 과제}
\label{sec:downstreams}

\paragraph*{신용 평가(Credit Scoring).}
사용 첫 12개월 안에 채무 불이행(default)이 발생할 확률을 예측해 리테일 신청에 대한 신용 위험(credit risk)을 평가하는 과제다.
다운스트림 데이터셋은 여러 해에 걸쳐 있으며 레코드가 다양하다.
이 과제는 소수 클래스(minority class)가 있는 이진 분류(binary classification) 문제로 정식화되며, 오프라인 지표로 ROC-AUC와 PR-AUC를 사용한다.

\paragraph*{커뮤니케이션 인게이지먼트(Communication Engagement).}
신용 신청을 도중에 중단한 사용자에게 발송된 재인게이지먼트(re-engagement) 커뮤니케이션을 사용자가 열지 여부를 예측하는 과제다.
이 행동은 신청 재개와 궁극적인 대출 실행으로 이어지는 깔때기(funnel)의 상단 대리 지표(proxy) 역할을 한다.
이 과제의 두드러진 점은 표본 크기가 매우 제한적이라는 것이며, 그래서 모델은 극히 적은 데이터에서 미묘한 이벤트 수준 신호를 잡아내야 한다.
이 과제는 이진 분류 문제로 정식화되며, 주된 오프라인 지표는 ROC-AUC와 PR-AUC이다.

\paragraph*{외부 사기(External Fraud).}
이는 이진 분류 문제로 정식화된 대표적인 사기 탐지 사용 사례다.
주된 오프라인 지표로 정밀도(precision)와 재현율(recall)을 사용해 성능을 평가한다.

\paragraph*{상품 추천(Product Recommendation).}
사용자가 특정 커뮤니케이션(예: 이메일 또는 푸시 알림)을 받은 조건에서 가까운 미래에 어떤 상품을 채택할 가능성이 있는지 예측하는 과제다.
핵심 난점은 커뮤니케이션의 맥락적 영향을 함께 고려하면서, 여러 상품에 대한 전환 성향(conversion propensity)을 동시에 모델링하는 것이다.
이 과제는 다중 라벨 분류(multilabel classification) 문제로 정식화되며, 모델은 포트폴리오 내 각 상품에 대해 독립적인 전환 확률을 출력한다.
성능은 평균 정밀도 평균(mean average precision, mAP)을 주요 오프라인 지표로 평가한다.

\paragraph*{정기 거래(Recurrent Transactions).}
주어진 거래가 다음 달에 반복될 정기 구독에 해당하는지 예측하는 과제다.
핵심 난점은 제한된 과거 신호 안에서 진짜 반복 패턴과 불규칙·일회성 결제를 구분해 내는 것이다.
이 문제는 이진 분류 과제로 정식화되며, 클래스 불균형을 반영하고 클래스 간 성능 균형을 맞추기 위해 매크로 평균 $F_\text{1}$-score로 성능을 평가한다.

\paragraph*{라이프타임 밸류(Lifetime Value, LTV).}
LTV 과제는 사용자가 양(陽)의 매출총이익(gross profit)을 발생시킬 확률을 평가하는 것으로, 이진 분류 문제로 정식화된다.
LTV 데이터셋의 두드러진 특징은 사용자들의 이벤트 이력이 짧다는 점이다(예: 몇 주). 반면 예측 지평(prediction horizon)은 보통 6개월 이상이다.
주된 오프라인 지표는 ROC-AUC와 PR-AUC이다.
